\chapter{Related Work}
\label{ch:related}

This thesis draws on five lines of work that are usually studied separately: parameter-efficient
fine-tuning, evidential and Bayesian uncertainty, calibration, out-of-distribution detection, and
few-shot learning. We review each with an eye on the design decision it informs, then position the
thesis against the closest prior work.

\section{Parameter-Efficient Fine-Tuning}

Houlsby et al.~\cite{houlsby2019adapters} introduced adapters for transformers: small bottleneck modules
(a down-projection, a non-linearity and an up-projection) inserted into a frozen network, with only the
bottlenecks trained. LoRA~\cite{hu2022lora} instead learns a low-rank update to an existing weight
matrix, $W = W_0 + BA$. Because the update can be merged into $W_0$ after training, LoRA adds no
inference latency. He et al.~\cite{he2022unified} show that adapters, LoRA and prefix tuning can be
written as variants of one update, differing in where the update is applied and whether it runs in
series with or in parallel to the frozen layer. We use their serial/parallel vocabulary. In vision, the
same ideas were carried over to Vision Transformers, for example visual prompt tuning~\cite{jia2022vpt}
and AdaptFormer~\cite{chen2022adaptformer}; a recent survey~\cite{xin2024peftsurvey} covers this body of
work and does not discuss calibration or OOD detection.

Work on CNNs is thinner. Conv-Adapter~\cite{chen2024convadapter} studies adapter design for ConvNets
and argues that adapters should respect the spatial structure of convolutional features, preferring
depth-wise $k \times k$ adapters to purely $1 \times 1$ ones. Convolutional variants of LoRA have been
proposed for domain adaptation~\cite{aleem2024convlora} and for the Segment Anything
model~\cite{zhong2024convlora}. Closest to our setting, Task-Specific Adapters (TSA)~\cite{li2022tsa}
attach adapters to a frozen ResNet-18 for cross-domain few-shot learning, classify with a
parameter-free nearest-centroid head, and find that parallel (residual) adapters outperform serial ones.
FiT~\cite{shysheya2023fit} adapts frozen CNN backbones with FiLM layers for few-shot transfer. For
vision-language models, CLIP-Adapter~\cite{gao2024clipadapter} and Tip-Adapter~\cite{zhang2022tipadapter}
attach small adapters to a frozen CLIP encoder. None of these works measures calibration and OOD
detection together with the adapter choice.

\section{Evidential and Bayesian Uncertainty}

A standard softmax classifier outputs one probability vector and has no way of saying ``none of the
above''. Bayesian methods place a distribution over weights, and deep
ensembles~\cite{lakshminarayanan2017ensembles} approximate this by training several networks, at a
multiplied training and inference cost. Evidential deep learning~\cite{sensoy2018edl} instead places a
distribution over the output. The network produces non-negative evidence $e_k$ for each class, which
defines a Dirichlet with parameters $\alpha_k = e_k + 1$, total strength $S = \sum_k \alpha_k$,
expected probabilities $\hat p_k = \alpha_k / S$ and vacuity $u = K/S$. Little evidence for any class
gives small $S$ and high vacuity, and all of this is available from a single deterministic forward
pass. Prior networks~\cite{malinin2018prior} use a related Dirichlet parameterisation but need OOD data
during training. Gao et al.~\cite{gao2024edlsurvey} survey the evidential literature, and Chen et
al.~\cite{chen2024redl} show that two standard ingredients of the EDL loss, the fixed ``$+1$'' prior and
the variance term, are not essential and can hurt calibration. Our evidential loss drops the variance
term accordingly (\cref{sec:edl-loss}).

The evidential approach has known weaknesses. Bengs et al.~\cite{bengs2022pitfalls} show that
uncertainty estimators trained by loss minimisation, standard EDL objectives included, need not behave
like faithful epistemic uncertainty, and Shen et al.~\cite{shen2024mirage} question how much of EDL's
apparent uncertainty quality is real. Hayta et al.~\cite{hayta2026plugin} show that the softmax
classifier is a special case of an evidential classifier under a family of plug-in losses, which helps
explain one of our RQ2 findings.

Bayesian ideas have also been combined with PEFT. Laplace-LoRA~\cite{yang2024laplacelora} fits a
post-hoc Laplace approximation over LoRA weights, and BLoB~\cite{wang2024blob} learns a variational
posterior over LoRA during training; both target large language models. In vision, Pandey et
al.~\cite{pandey2024bayesianpeft} find that PEFT-adapted vision foundation models are underconfident in
few-shot settings and propose a Bayesian PEFT framework with base-rate adjustment and an evidential
ensemble. BayesAdapter~\cite{moralesalvarez2026bayesadapter} learns a probability distribution over the
parameters of a CLIP adapter instead of a point estimate and reports better calibration and selective
classification. Bayesian Evidential Learning
(BEL)~\cite{linghu2022bel} brings evidential uncertainty to few-shot classification; on miniImageNet
5-shot it lowers the ECE of a Meta-Baseline model with a Conv-4 backbone from 14.69\% to 3.59\%. Its
backbone is meta-trained end to end, unlike ours.

\section{Calibration}

Guo et al.~\cite{guo2017calibration} showed that modern deep networks are overconfident, popularised
the expected calibration error (ECE) as a summary measure, and proposed temperature scaling: dividing
the logits by one scalar fitted on held-out data by minimising the negative log-likelihood. Because it
is a monotone rescaling, temperature scaling never changes the predicted class. Minderer et
al.~\cite{minderer2021revisiting} revisit calibration for newer architectures and find that model family
matters a great deal, and Tao et al.~\cite{tao2024benchmark} study calibration across a large space of
architectures. Neither decomposes calibration variance across design factors the way RQ1 does. Dirichlet
calibration~\cite{kull2019dirichlet} uses the word ``Dirichlet'' in a different sense: it is a post-hoc
calibration map for softmax outputs and is unrelated to evidential Dirichlet evidence. More flexible
accuracy-preserving calibrators also exist for softmax classifiers~\cite{zhao2026invlt}. For evidential
models, recent fixes change the model or its training, for example density-informed
pseudo-counts~\cite{carlotti2026dipedl} or a self-calibrated evidence-strength head~\cite{turacan2026invascal};
neither is a small post-hoc refit of the kind studied in RQ4.

\section{Out-of-Distribution Detection}

Hendrycks and Gimpel~\cite{hendrycks2017baseline} proposed the maximum softmax probability (MSP) as a
baseline OOD score and established AUROC as the standard evaluation. Liu et al.~\cite{liu2020energy}
proposed the energy score, the log-sum-exp of the logits, which needs no retraining and separates OOD
inputs better than MSP. OpenOOD~\cite{yang2022openood,zhang2024openood15} standardised the distinction
between far-OOD (semantically unrelated data, such as digits for an animal classifier) and near-OOD
(related but unseen classes) that we follow. McNamara~\cite{mcnamara2026vacuity} shows that
vacuity-based AUROC can be inflated when the in-distribution and OOD evaluations use different numbers
of classes; in our protocol every OOD image is scored by the same 5-way episode head, so this artefact
does not arise.

Two recent studies are close to RQ2. Gen{\c{c}} et al.~\cite{genc2026objectives} compare training
objectives for OOD detection and note that, in prior work, each objective tends to be evaluated with
the confidence measure natural to its output space, so objective and score are not fully separated.
Their objectives do not include an evidential one, and they do not study frozen backbones or episodic
evaluation. Krumpl et al.~\cite{krumpl2026onemodel} cross 56 ImageNet-trained ResNet-50 models with 21
post-hoc detectors on 8 OOD sets and find a strong dependence between training strategy and detector.
We run a smaller crossing of the same kind, with evidential models included.

\section{Few-Shot Learning}

Few-shot classification is usually evaluated with $N$-way $K$-shot episodes~\cite{vinyals2016matching}:
each episode samples $N$ classes, $K$ labelled support images per class, and a set of query images to
classify. We use the MiniImageNet split of Ravi and Larochelle~\cite{ravi2017optimization} and the
CIFAR-FS split of Bertinetto et al.~\cite{bertinetto2019r2d2}. Prototypical networks~\cite{snell2017protonet}
classify a query by its distance to each class prototype, the mean embedding of the class's support
images, and need no trainable classifier. Chen et al.~\cite{chen2019closer} show that simple baselines
with a cosine-similarity classifier are strong few-shot competitors. MetaOptNet~\cite{lee2019metaoptnet}
is a representative method trained from scratch on the base classes, while P>M>F~\cite{hu2022pmf} shows
how much large-scale pretraining (for example DINO~\cite{caron2021dino}) adds when the whole backbone is
then meta-trained and fine-tuned.

\section{Backbones for Edge Deployment}

We use ResNet-18~\cite{he2016resnet} and MobileNetV3-Small~\cite{howard2019mobilenetv3}, both
pretrained on ImageNet. They were chosen because they are realistic for constrained hardware, not for
peak accuracy. A recent survey of tiny machine learning~\cite{somvanshi2025tinydl} describes the memory
and compute limits of such hardware and the continued role of compact CNNs there.

\section{Positioning}

\Cref{tab:gap} compares this thesis with the three closest pieces of prior work. Each covers some of
the four properties we study, but none covers all of them. In addition, we separate two pairs of factors
that published comparisons usually change together: the training objective versus the OOD scoring rule
(RQ2), and adapter architecture versus parameter budget (RQ3).

\begin{table}[H]
\centering
\caption{The research gap. Each close prior work covers some, but not all, of the four properties
studied here.}
\label{tab:gap}
\begin{adjustbox}{max width=\textwidth}
\begin{tabular}{@{}lcccc@{}}
\toprule
& Frozen lightweight CNN & Parameter-efficient adapter & Calibration & OOD detection \\
\midrule
TSA~\cite{li2022tsa} & \cmark & \cmark & \xmark & \xmark \\
BEL~\cite{linghu2022bel} & \xmark\ (meta-trained) & \xmark & \cmark & \xmark \\
BayesAdapter~\cite{moralesalvarez2026bayesadapter} & \xmark\ (frozen CLIP) & \cmark & \cmark & \xmark \\
This work & \cmark & \cmark & \cmark & \cmark \\
\bottomrule
\end{tabular}
\end{adjustbox}
\end{table}
